\chapter{Meromorphic Functions}

\section{Zeros and Poles}

\begin{definition}{Isolated Singularities}\\
An \textit{isolated singularity} (or \textit{point singularity}) of a function $f$ is a complex number $z_0$ such that $f$ is defined in an open neighborhood $z_0$ except for the point $z_0$ itself.
\end{definition}

\noindent\textbf{Motivation:} The singularities of holomorphic functions usually appear when the denominator in a fraction equals to zero, so to fully master singularities, we need to study the zeros of functions first.

\begin{definition}{Zeros}\\
A complex number $z_0$ is a \textit{zero} for the holomorphic function $f$ if $f(z_0)=0$.
\end{definition}

\noindent\textbf{Remark:} The identity Thm implies that the zeros of a non-trivial holomorphic function are isolated.

\begin{theorem}{Local Factorization Theorem for Zeros}\\
Let $f$ be holomorphic in a connected open subset $\Omega \subseteq \mathbb{C}$. Suppose $f$ has a zero at a point $z_0 \in \Omega$ and doesn't vanish identically in $\Omega$. Then there is an open neighborhood $U\subseteq \Omega$ of $z_0$, a non-vanishing holomorphic function $g$ on $U$, and a unique positive integer $n$, s.t.: $\forall z\in U$,
\[
f(z)= (z-z_0)^n g(z)
\]
\end{theorem}

\begin{proof}
    In a small disc centered at $z_0$, $f$ has a power series expansion $f(z)= \sum_{k=0}^\infty a_k(z-z_0)^k$. Since $f$ is not identically zero, then there is a smallest positive integer $n$, s.t.: $a_n\neq 0$. Then we write
    \[
    f(z)=(z-z_0)^n (a_n+ a_{n+1}(z-z_0)+ \cdots)
    \]
    Define $g(z):= \sum_{k=9}^\infty a_{n+k} (z-z_0)^k$, which is holomorphic and is non-vanishing $\forall z$ closed to $z_0$ (since $a_n \neq 0$). To see the uniqueness of $n$, we suppose that $f(z)= (z-z_0)^n g(z)= (z-z_0)^m h(z)$ for some $m\in \mathbb{N}_+$ and some $h(z_0) \neq 0$. Suppose on the contrary that, WLOG, $m>n$. Then $g(z)= (z-z_0)^{m-n} h(z)$, where $z\to z_0$ yields $g(z_0)=0$, contradiction. Hence, $m=n$.
\end{proof}

\begin{definition}{Order of Zeros}\\
The unique $n\in \mathbb{N}_+$ in the Local Factorization Theorem for Zeros is called the \textit{order} (or \textit{multiplicity}) of the zero $z_0$ for $f$. Furthermore, if $n=1$, then $z_0$ is said to be \textit{simple}.
\end{definition}

\begin{definition}{Poles}
\begin{itemize}
    \item A \textit{deleted neighborhood} of $z_0$ is defined to be an open disc
    \[
    \{z\mid 0<|z-z_0|<r\}
    \]
    for some $r>0$.

    \item A function $f$ defined in a deleted neighborhood of $z_0$ has a \textit{pole} at $z_0$ if the function $1/f$, defined to be zero at $z_0$, is holomorphic in a full neighborhood of $z_0$.
\end{itemize}
\end{definition}

\begin{theorem}{Local Representations of Poles Theorem}\\
If $f$ has a pole at $z_0 \in \Omega$, then in an open neighborhood of $z_0$, there exist a non-vanishing holomorphic function $h$ and a unique positive integer $n$, s.t.:
\[
f(z)= (z-z_0)^{-n} h(z)
\]
\end{theorem}

\begin{proof}
    It's immediately from the def of poles and the Local Factorization Thm for Zeros.
\end{proof}

\begin{definition}{Order of Poles}\\
The unique $n$ in the Local Representations of Poles Thm is called the \textit{order} (or \textit{multiplicity}) of the pole $z_0$ for $f$. Furthermore, if $n=1$, then $z_0$ is said to be \textit{simple}.
\end{definition}

\begin{theorem}{Principal Part Expansions at Poles Theorem}\\
If $f$ has a pole of order $n$ at $z_0$, then
\[
f(z)= \frac{a_{-n}}{(z-z_0)^n} +\frac{a_{-n+1}}{(z-z_0)^{n-1}} +\cdots+ \frac{a_{-1}}{z-z_0} +G(z)
\]
where $G$ is a holomorphic function in an open neighborhood of $z_0$.
\end{theorem}

\begin{proof}
    It's immediately from the Local Representations of Poles Thm by expanding $h$.
\end{proof}

\begin{definition}{Principal Part}\\
The sum
\[
\frac{a_{-n}}{(z-z_0)^n} +\frac{a_{-n+1}}{(z-z_0)^{n-1}} +\cdots+ \frac{a_{-1}}{z-z_0}
\]
in the Principal Part Expansions at Poles Thm is called the \textit{principal part} of $f$ at the pole $z_0$.
\end{definition}

\begin{definition}{Residue}\\
The coefficient $a_{-1}$ in the Principal Part Expansions at Poles Thm is the \textit{residue} of $f$ at the pole, denoted as $\text{res}_{z_0} f:=a_{-1}$.
\end{definition}

\noindent\textbf{Remark:} \begin{itemize}
    \item $\forall k=2,\cdots, n$, the term $\frac{a_{-k}}{(z-z_0)^{k}}$ in the principal part has primitive in a deleted neighborhood of $z_0$:
    \[
    \frac{-a_{-k}}{(k-1)(z-z_0)^{k-1}}
    \]
    where integral on any circle $C$ centered at the origin is zero and which is a single-valued function. But for the residue $a_{-1}$ of $f$, its primitive is $\text{log } (z-z_0)$, which is multi-valued. This is why we are discussing the situation of $a_{-1}$ separately.

    \item If $f$ has a simple pole at $z_0$, then it's clear that
    \[
    \text{res}_{z_0}f= \lim_{z\to z_0} (z-z_0)f(z)
    \]
    But for the pole of higher order, the form of evaluation of $\text{res}_{z_0}f$ is not trivial:
\end{itemize}

\begin{theorem}{Residue Computation Theorem}\\
If $f$ has a pole of order $n$ at $z_0$, then
\[
\text{res}_{z_0}f= \lim_{z\to z_0} \frac{1}{(n-1)!} (\frac{d}{dz})^{n-1} (z-z_0)^n f(z)
\]
\end{theorem}

\begin{proof}
    It's immediately from a consequence of the Principal Part Expansions at Poles Thm:
    \[
    (z-z_0)^n f(z)= a_{-n} +a_{-n+1} (z-z_0) +\cdots+ a_{-1} (z-z_0)^{n-1}+ G(z)(z-z_0)^n
    \]
\end{proof}

\section{Residue Formula}

\begin{theorem}{Cauchy's Residue Theorem}\\
Suppose $f$ is holomorphic in an open set containing a circle $C$ and its interior, except for a pole at $z_0$ inside $C$. Then
\[
\int_C f(z)dz= 2\pi i\text{res}_{z_0}f
\]
\end{theorem}

\begin{proof}
    Consider a keyhole contour whose small circle $C_\epsilon$ is centered at $z_0$ and of radius $\epsilon$. Let the width of the corridor go to zero and then we have:
    \[
    \int_C f(z)dz= \int_{C_\epsilon} f(z)dz
    \]
    Note that by applying CIF to the constant function $a_{-1}$, we get:
    \[
    \frac{1}{2\pi i} \int_{C_\epsilon} \frac{a_{-1}}{z-z_0}dz=a_{-1} =\text{res}_{z_0}f
    \]
    that $\forall k>1$, by applying coro.1 of CIF to the constant function $a_{-k}$, we get:
    \[
    \frac{1}{2\pi i} \int_{C_\epsilon} \frac{a_{-k}}{(z-z_0)^k} dk=0
    \]
    and that by Cauchy's Thm, we have: $\int_{C_\epsilon} G(z)dz=0$, since $G$ is holomorphic. Hence, applying $\int_{C_\epsilon} * dz$ to $f(z)$ with
    \[
    f(z)= \frac{a_{-n}}{(z-z_0)^n} +\frac{a_{-n+1}}{(z-z_0)^{n-1}} +\cdots +\frac{a_{-1}}{z-z_0} +G(z)
    \]
    we're done.
\end{proof}

\begin{corollary}{1 of Cauchy's Residue Theorem}\\
Suppose $f$ is holomorphic in an open set containing a circle $C$ and its interior, except for poles at the points $z_1 ,\cdots, z_N$ inside $C$. Then
\[
\int_Cf(z)dz= 2\pi i\sum_{k=1}^N \text{res}_{z_k}f
\]
\end{corollary}

\begin{proof}
    Similar as the proof of Cauchy's Residue Thm, except for considering a multiple keyhole contour to passby each pole.
\end{proof}


\begin{corollary}{2 of Cachy's Residue Theorem: Residue Formula}\\
Suppose $f$ is holomorphic in an open set containing a toy contour $\gamma$ and its interior, except for poles at the points $z_1 ,\cdots, z_N$ inside $\gamma$. Then
\[
\int_\gamma f(z)dz= 2\pi i\sum_{k=1}^N \text{res}_{z_k}f
\]
\end{corollary}

\section{Singularities}

\noindent\textbf{Intuition:} We've learnt something  about the poles, so now we need to pay attention to some other types of isolated singularities.

\subsection{Removable Singularities}

\noindent\textbf{Motivation:} Consider the function $\frac{(z-1)(z-2)}{z-2}$, which has a singularity at $z=2$, but can be extended to be $z-1$.

\begin{definition}{Removable Singularities}\\
Let $f$ be a function holomorphic in an open set $\Omega \subseteq \mathbb{C}$, except for the point $z_0 \in \Omega$. We say that $z_0$ is a \textit{removable singularity} if we can define $f$ at $z_0$ in such a way that $f$ becomes holomorphic in all of $\Omega$.
\end{definition}

\begin{theorem}{Riemann's Theorem on Removable Singularities}\\
Suppose $f$ is holomorphic in an open set $\Omega$, except for a point $z_0 \in \Omega$. If $f$ is bounded on $\Omega \setminus \{z_0\}$, then $z_0$ is a removable singularity.
\end{theorem}

\begin{proof}
    Consider an open disc $D$ centered at $z_0$ whose closure is contained in $\Omega$, let $C$ denote the boundary of $D$, and fix $z\in D$ with $z\neq z_0$. Consider a multiple keyhole contour with two small circles where $z_0$ and $z$ are the centers of two circles $C_0$ respectively $C_1$ both of radius $\epsilon$. Let the sides of the two corridors go close to each other and finally overlap. Then we get:
    \[
    \int_C \frac{f(\zeta)}{\zeta -z}d\zeta +\int_{C_0} \frac{f(\zeta)}{\zeta -z}d\zeta + \int_{C_1} \frac{f(\zeta)}{\zeta -z}d\zeta=0
    \]
    As in the proof of CIF, we have:
    \[
    \int_{C_1} \frac{f(\zeta)}{\zeta -z}d\zeta=-2\pi if(z)
    \]
    Take a very small $\epsilon$ such that $\zeta$ on $C_0$ stays from $z$. Then by the assumption that $f$ is bounded, i.e.: $|f|\leq M$ for some $M>0$, we have:
    \[
    |\int_{C_0} \frac{f(\zeta)}{\zeta -z}d\zeta| \overset{(\text{ML)}}{\leq} \frac{M}{\sup_{\zeta \in C_0} |\zeta -z|} 2\pi \epsilon \xrightarrow{\epsilon\to0}
    \]
    Therefore, we can conclude that
    \[
    f(z)= \frac{1}{2\pi i} \int_{C} \frac{f(\zeta)}{\zeta -z}d\zeta
    \]
    Then by the coro of Identity Thm, we're done.
\end{proof}

\begin{corollary}{of Riemann's Theorem on Removable Singularities: Equivalent Characterization of Poles Proposition}\\
Suppose $f$ has an isolated singularity at the point $z_0$. Then $z_0$ is a pole of $f$ iff $|f(z)| \to \infty$ as $z\to z_0$.
\end{corollary}

\begin{proof}
    $(\Rightarrow)$ holds immediately by the def of poles. Conversely, by assumption, $1/f(z)$ is bounded near $z_0$. Hence, $1/f$ has a removable singularity at $z_0$ and in fact vanish at $z_0$. Therefore, by the def of removable singularities, we can extend $1/f$ to be zero at $z_0$, so $f$ has a pole at $z$.
\end{proof}

\subsection{Essential Singularities}

\begin{definition}{Essential Singularities}\\
Any singularity that is not removable or a pole is defined to be an \textit{essential singularity}.
\end{definition}

\noindent\textbf{Example:} $e^{1/z}$ has an essential singularity at $z=0$: Near $z=0$, $e^{1/z}$ oscillates and grows faster than any power, and a complete understanding of its behavior is not easy.

\begin{theorem}{Casorati-Weierstrass Theorem}\\
Suppose $f$ is holomorphic in the punctured disc $D_r(z_0) \setminus \{z_0\}$ and has an essential singularity at $z_0$. Then the image of $D_r(z_0) \setminus \{z_0\}$ under $f$ is dense in the complex plane.
\end{theorem}

\begin{proof}
    Suppose on the contrary that $\exists w\in \mathbb{C}, \exists \delta>0,$ s.t.: $\forall z\in D_r(z_0) \setminus \{z_0\}, |f(z)-w|>\delta$. Consider the function
    \[
    g(z):= \frac{1}{f(z)-w}
    \]
    with $|g(z)|< \frac{1}{\delta}$ and hence bounded. Then by the Riemann's Thm on Removable Singularities, $g(z)$ has a removable singularity at $z_0$:
    \begin{itemize}
        \item If $g(z_0) \neq 0$, then $f(z)= w+\frac{1}{g(z)}$ is holomorphic at $z_0$, contradiction.
        \item If $g(z_0)=0$, then $f(z)-w= \frac{1}{g(z)}$ has a pole at $z_0$, contradiction.
    \end{itemize}
\end{proof}

\section{Meromorphic Functions}

\begin{definition}{Meromorphic Functions}\\
A function $f$ on an open set $\Omega$ is \textit{meromorphic} if there exists a sequence of points $\{z_0,z_1 ,\cdots\}$ that has no limit points in $\Omega$ satisfying the followings:
\begin{itemize}
    \item $f$ is holomorphic in $\Omega \setminus \{z_0 ,z_1 ,\cdots\}$.
    \item $f$ has poles at the points $\{z_0 ,z_1,\cdots\}$.
\end{itemize}
\end{definition}

\noindent\textbf{Remark} (\textit{Meromorphic Functions in the Extended Complex Plane}). If a function $f$ is holomorphic for all large values of $z$, then the function $F(z):= f(1/z)$ is holomorphic in a deleted neighborhood of the origin.
\begin{itemize}
    \item We then say that $f$ has a \textit{pole at $\infty$} if $F$ has a pole at the origin. Similarly, we can speak of $f$ having an \textit{essential singularity at $\infty$} or a \textit{removable singularity} (hence \textit{holomorphic}) \textit{at $\infty$} in terms of the corresponding behavior of $F$ at $0$.

    \item A meromorphic function in $\mathbb{C}$, that is either holomorphic at $\infty$ or has a pole at $\infty$, is said to be a meromorphic function in the extended complex plane $\widehat{\mathbb{C}}:= \mathbb{C} \cup \{\infty\}$.
\end{itemize}

\begin{theorem}{Characterization of Meromorphic Functions in the Extended Complex Plane Theorem}\\
The meromorphic functions in $\widehat{\mathbb{C}}$ are the rational function.
\end{theorem}

\begin{proof}
    Let $f$ be meromorphic in $\widehat{\mathbb{C}}$. Then $f$ has finitely many poles $z_1 ,\cdots, z_n$ in $\widehat{\mathbb{C}}$, since otherwise the sequence of infinitely many points in $\widehat{\mathbb{C}}$ must converge to somewhere in $\widehat{\mathbb{C}}$, which is a contradiction to the def of meromorphic functions. $\forall k=1 ,\cdots, n,$ near $z_k$, let $f_k(z)$ denote the principal part of $f(z)$, and at $\infty$, let $\widetilde{f}_\infty (z)$ denote the principal part of $f(1/z)$ at $0$ and then set $f_\infty (z):= \widetilde{f}_\infty (1/z)$. Consider a function 
    \[
    H(z):= f(z)- f_\infty (z) -\sum_{k=1}^n f_k(z)
    \]
    which is entire and bounded in all of $\widehat{\mathbb C}$. By Liouville's, $H(z) \equiv M$ is a contant function. Hence,
    \[
    f(z)+ f_\infty (z)+ \sum_{k=1}^n f_k(z) +M
    \]
    Note the every principal part $f_k$ is polynomial in $1/(z-z_k)$ and that $\widetilde{f} _\infty$ is a polynomial in $1/z$ which implies that $f_\infty$ is a polynomial in $z$. Therefore, $f$ is a rational function.
\end{proof}

\section{Argument Principle}

\subsection{Argument Principle}

\noindent\textbf{Intuition:} Recall that the function $\text{log }f(z) := \text{log }|f(z)| +i\text{arg }f(z)$ is multi-valued, but the derivative is $f'(z)/f(z)$, which is single-valued. So we are interested in studying $f'/f$.

\begin{theorem}{Argument Principle}\\
Suppose $f$ is meromorphic in an open set containing a circle $C$ and its interior. If $f$ has no poles and never vanishes on $C$, then
\[
\frac{1}{2\pi i} \int_C \frac{f'(z)}{f(z)}dz= \# (\text{zeros of }f \text{ inside }C)- \# (\text{poles of }f \text{ inside }C)
\]
where the zeros and poles are counted with their multiplicities.
\end{theorem}

\begin{proof}
    \begin{itemize}
        \item If $f$ is holomorphic and has a zero of order $n$ at $z_0$, we write
        \[
        f(z)= (z-z_0)^n g(z)
        \]
        where $g$ is holomorphic and non-vanishing in a neighborhood of $z_0$. Hence, 
        \[
        \frac{f'(z)}{f(z)}= \frac{n}{z-z_0} +\frac{g'(z)}{g(z)}
        \]
        has a simple pole with residue $n$ at $z_0$.

        \item If $f$ is holomorphic and has a pole of order $n$ at $z_0$, then we write $f(z)= (z-z_0)^{-n} h(z)$ and then we have
        \[
        \frac{f'(z)}{f(z)} =\frac{-n}{z-z_0} +\frac{h'(z)}{h(z)}
        \]
        having a simple pole with residue $-n$ at $z_0$.
    \end{itemize}
    Therefore, if $f$ is meromorphic, then $f'/f$ have simple poles at the zeros and poles of $f$, with residue $=$ the order of the zero of $f$ or the negative of the order of the pole of $f$.
\end{proof}

\noindent\textbf{Remark:} This thm holds also for toy contours.

\subsection{Applications of Argument Principle}

\subsubsection{Rouche's Theorem}

\begin{theorem}{Rouche's Theorem}\\
Suppose $f$ and $g$ are holomorphic in an open set containing a circle $C$ and its interior. If $\forall z\in C$,
\[
|f(z)| >|g(z)|
\]
then $f$ and $f+g$ have the same number of zeros inside $C$.
\end{theorem}

\begin{proof}
    For $t\in [0,1]$, define $f_t(z) := f(z) +tg(z)$. Then $f_0=f$ and $f_1=f+g$. Let $n_t$ denote the number of zeros of $f_k$ inside $C$ counted with multiplicities. $\forall z\in C$, since $|f(z)| >|g(z)|$, then
    \[
    |f_t(z)|\geq |f(z)|- t |g(z)| \geq |f(z)|- |g(z)|>0
    \]
    Hence, $f_t$ has no zeros on $C$, so by the Argument Principle, we have
    \[
    n_t= \frac{1}{2\pi i} \int_C \frac{f'_t(z)}{f_t(z)}dz
    \]
    Therefore, $n_t$ is continuous for $t\in [0,1]$. By the def of $n_t$, we then see that it must be a constant.
\end{proof}

\subsubsection{Open Mapping Theorem}

\begin{theorem}{Open Mapping Theorem}\\
If $f$ is holomorphic and non-constant in an open set $\Omega \subseteq \mathbb C$, then $f$ is open.
\end{theorem}

\begin{proof}
    $\forall w_0 \in \text{im }f, \exists z _0\in \Omega$, s.t.: $f(z_0)=w_0$. Define 
    $$g(z) := f(z)-w, F(z):= f(z)-w_0 \text{ and } G(z):= w_0-w$$
    Then $g=F+G$. Since $\Omega$ is open and $f$ is not constant, then $\exists \delta>0$, s.t.: the closure of the disc $D: |z-z_0|<\delta$ is contained in $\Omega$ and that $f(z) \neq w_0$ on the boundary $C$ of $D$. Then we take $\epsilon>0$ to be such that $|f(z)-w_0|\geq \epsilon$ on $C$. Consider $B_\epsilon (w_0)$. $\forall w\in B_\epsilon (w_0), |G(z)|< \epsilon \leq F|(z)|$ on $C$, so can apply Rouche's Thm to conclude that $F$ has the same number of zeros of $g=F+G$ inside $C$, say $g$ has at least one zero inside $C$, i.e.: $\forall w\in B_\epsilon (w_0), \exists z\in D\subseteq \Omega$, s.t.:
    \[
    f(z)=w
    \]
    which implies that $B_\epsilon (w) \subseteq \text{im }f$ and hence gives the openness of $\text{im }f$.
\end{proof}

\subsubsection{Maximum Modulus Principle}

\begin{definition}{Maximum of Holomorphic Function}\\
The \textit{maximum} of a holomorphic function $f$ in an open set $\Omega$ is defined to be the maximum of its absolute value $|f|$ in $\Omega$.
\end{definition}

\begin{theorem}{Maximum Modulus Principle}\\
If $f$ is a non-constant holomorphic function in an open set $\Omega  \subseteq \mathbb{C}$, then $f$ cannot attain a maximum in $\Omega$.
\end{theorem}

\begin{proof}
    Suppose on the contrary that $f$ attain a maximum ta $z_0 \in \Omega$. By the Open Mapping Thm, $f$ is open. Hence, $f(D)$ is open and is containing $f(z_0)$, where $D\subseteq \Omega$ is an arbitrary open disc containing $z_0$. Since $f$ is non-constant, then we can find some $D$, s.t.: $\exists z\in D$ woth $|f(z)|> |f(z_0)|$, contradiction.
\end{proof}

\begin{corollary}{of Maximal Modulus Principle}\\
Suppose that $\Omega$ is an open set with compact closure. If $f$ is holomorphic on $\Omega$ and continuous on $\overline{\Omega}$, then
\[
\sup_{z\in \Omega} |f(z)| \leq \sup_{z\in \overline{\Omega} \setminus \Omega} |f(z)|
\]
\end{corollary}